
\subsection*{5.2}
\phantomsection
\label{sga3:VIIB:5.2}

Throughout the rest of \SGARefLink{sga3:VIIB:5}{Section~5}, \(k\) denotes a perfect field
of characteristic \(p>0\).

Put
\[
  \overline{\mathbb N}=\mathbb N\cup\{\infty\}.
\]
If \(B\) is a profinite \(k\)-algebra and
\(r\in\overline{\mathbb N}\), denote by
\[
  \bigl((x^{p^r})\bigr)_{x\in\mathfrak r(B)}
\]
the closed ideal of \(B\) generated by the elements \(x^{p^r}\),
where \(x\) ranges over the radical \(\mathfrak r(B)\) of \(B\).
If \(r=\infty\), use the same notation with the convention that
\[
  \bigl((x^{p^\infty})\bigr)_{x\in\mathfrak r(B)}
\]
is the zero ideal. In both cases, \(B_r\) denotes the quotient
\[
  B_r
  =
  B\big/\bigl((x^{p^r})\bigr)_{x\in\mathfrak r(B)}.
\]

We say that \(B\) is \emph{of height at most \(r\)} if
\[
  \bigl((x^{p^r})\bigr)_{x\in\mathfrak r(B)}
\]
is the zero ideal. If this holds and \(r\) is finite, we say that
\(B\) is \emph{of finite height}.

Consider in particular the case
\[
  B=k[[\omega]],
\]
where \(\omega\) is a pseudocompact \(k\)-vector space
\textup{(cf.\ \SGARefLink{sga3:VIIB:1.2.5}{1.2.5})}.\footnote{%
\label{sga3:VIIB:nde:153}In the original, the author uses
``linearly compact vector space,'' which is equivalent to
``pseudocompact vector space'' \textup{(cf.\ [BAC],
\S III.2, Exercises 15\,a), 19\,a), and 20\,d))}. The editors
preferred to retain the terminology ``pseudocompact,'' used up to
this point.}
We then say that \(B\) is a \emph{formal power series algebra}
and that \(B_r\) is a \emph{truncated formal power series algebra}
\(\bigl(r\in\overline{\mathbb N}\bigr)\); thus, by convention,
\[
  B=B_\infty
\]
is also called ``truncated.'' If \(B=k[[\omega]]\), we also write
\[
  \bigl((x^{p^r})\bigr)_{x\in\omega}
\]
instead of
\[
  \bigl((x^{p^r})\bigr)_{x\in\mathfrak r(B)}.
\]

\medskip
\noindent\textbf{Notation.}
Let \(\omega\) be a pseudocompact \(k\)-vector space filtered by
an increasing sequence of closed vector subspaces
\[
  0=\omega_0\subset\omega_1\subset\omega_2
  \subset\omega_3\subset\cdots.
\]
\begin{enumerate}
\item[\textup{\SGAInlineRefTarget{sga3:VIIB:5.2:item:a}(a)}]
The closed ideal of \(k[[\omega]]\) generated by the elements
\(x^{p^r}\), where \(r\) ranges over \(\mathbb N\) and
\(x\) ranges over \(\omega_r\), will be denoted by
\[
  \bigl((x^{p^r})\bigr)_{x\in\omega_r}.
\]

\item[\textup{(b)}]
For \(r\in\mathbb N\), denote by \({}^r\omega\) the filtered
pseudocompact vector space defined by
\[
  {}^r\omega_i=\omega_i\quad\text{if }i<r,
  \qquad
  {}^r\omega_i=\omega\quad\text{if }i\geq r.
\]
\end{enumerate}

\phantomsection
\label{sga3:VIIB:5.2:notation-source-note}
\noindent\textit{Source note.}\enspace
In Notation \textup{\SGARefLink{sga3:VIIB:5.2:item:a}{(a)}}, the French PDF prints the subscript
\(r\in\omega_r\). The preceding defining sentence and every
subsequent use require \(x\in\omega_r\), which is used here.

\medskip
\noindent\textbf{Theorem (Dieudonn\'e--Cartier).}
\phantomsection
\label{sga3:VIIB:5.2:theorem}
{\itshape Let
\[
  H\longrightarrow G
\]
be a monomorphism of infinitesimal formal groups over a perfect
field \(k\) of characteristic \(p>0\). Let \(B\) be the affine
algebra of the homogeneous space \(G/H\), and suppose that one of
the following three conditions holds:\footnote{%
\label{sga3:VIIB:nde:154}The editors replaced \(H\backslash G\)
by \(G/H\), and did the same in the proof; they also added
condition \textup{\SGARefLink{sga3:VIIB:5.2:theorem:item:iii}{(iii)}}.}
\begin{enumerate}
\item[\textup{(i)}]
\(B\) has finite height \textup{(in particular, this holds if
\(G\) has finite height)}.

\item[\textup{\SGAInlineRefTarget{sga3:VIIB:5.2:theorem:item:ii}(ii)}]
\(B\) is a complete Noetherian local ring.

\item[\textup{\SGAInlineRefTarget{sga3:VIIB:5.2:theorem:item:iii}(iii)}]
\(B\) is a reduced ring.
\end{enumerate}
Then \(B\) is isomorphic to the completed tensor product of a
finite family of truncated formal power series algebras.}

The proof of this theorem occupies \SGARefLink{sga3:VIIB:5.2.1:section-proposition}{5.2.1}--\SGARefLink{sga3:VIIB:5.2.5}{5.2.5}.
